% Nouveau EGA V (cote 25): archive cote 25, site page 161 (PDF page 162).
% typescript. Draft compared: none.
% AI-transcribed by Claude Opus 5.5 (claude-opus-5-5, adaptive thinking, effort high), Message Batches API, batch ega25b (msgbatch_01A7n5YuXubUPNK1zsepi8RG), 2026-11-13; not yet checked by eye.
\scanpage{25/161}
\mspage{\hand{\unsure{4|H}}}

puissance de chaque v\unsure{a|X}riable $X_i$ . Plus précis\var{\xout{\unsure{m}}}{}{typewriter strike-out}{ega25_25_161-v0001}éme\unsure{n|t}t, le sous-\var{\xout{pré}}{}{typewriter strike-out}{ega25_25_161-v0002}schéma
de $P^r$ défini par l'idéal homogène engendré par \var{\xout{le}}{}{typewriter strike-out}{ega25_25_161-v0003} $\hand{\phi}$ et les $\hand{\frac{\partial\hand{\phi}}{\partial\textcolor{black}{X_i}}}$ est
exactement formé des points de V($\hand{\phi}$) en lesquels V($\hand{\phi}$) n'est pas lisse sur
Spec(A) avec une dimension relative r-1.
15.2. Soi\unsure{e}nt X un préschéma, $\underline{J}$ un Idéal \var{\inserted[typed]{qu coh}}{qu coh}{insertion}{ega25_25_161-v0004} \var{\xout{de \unsure{ty}}}{}{typewriter strike-out}{ega25_25_161-v0005} sur X, nous dirons que
l'homomorphisme ``d'augmentation'' $\underline{O}_X \hand{\to} \underline{O}_X/\underline{J}$ est une ``augmentation
élémentaire de multiplicité n'' si elle satisfait aux conditions suivantes
a) $\underline{J}/\underline{J}^2 = \underline{N}_{Y/X}$ est loc libre de type fini sur Y = V($\underline{J}$).
b) Le noyau de l'homomorphisme canonique
\[ \var{\over{\unsure{u}}{\mu}}{\mu}{writing over}{ega25_25_161-v0006} : \underline{Sym}_{\underline{O}_Y}(\underline{N}) \hand{\longrightarrow} gr_{\underline{J}}(\underline{O}_X) \]
est engendré par le no\unsure{y|g}\unsure{a|g}u \var{\inserted[typed]{K}}{K}{insertion}{ega25_25_161-v0007} de $\var{\over{\unsure{u}}{\mu}}{\mu}{writing over}{ega25_25_161-v0008}^n$, lequel est un sous-Module inversible
localement facteur direct de $\underline{Sym}^n_{\underline{O}_Y}(\underline{N})$ (i.e. est engendrable localement
sur Y par une section $\hand{\phi}$ qui ne\var{\xout{est}}{}{typewriter strike-out}{ega25_25_161-v0009} s'annule en aucun point).
c) Ladite $\hand{\phi}$ (qui n'est définie, localement, que modulo multiplication
par une unité) est une forme \uline{lisse} . (NB On aurait pu introduire dans
\var{\margin{\textbar}}{\placedmargin{\textbar}}{marginal note placed}{ega25_25_161-v0010}15.1. la notion de sous-Module inversible lisse de $\underline{Sym}^n(E)$, à vrai dire
géométriquement plus importante que celle de \uline{section} lisse, car V($\hand{\phi}$) est
en fait défini par un tel sous-Module ...\var{\xout{\unsure{;}}}{}{typewriter strike-out}{ega25_25_161-v0011} cela aurait l'avantage de
permettre de formuler b) et c) en une seule condition).

Lorsque n=2, on parlera de ``augmentation \var{\inserted[typed]{de type}}{de type}{insertion}{ega25_25_161-v0012} quadratique élémentaire'' ,
Si A est un anneau, \var{\xout{local noethérien}}{}{typewriter strike-out}{ega25_25_161-v0013} muni d'une \unsure{augmentation|augmantation} A$\hand{\to}$A/J,
on conviendra encore de dire que cette augmentation est élémentaire de
multiplicité \unsure{n} , s'il en \unsure{est} ainsi pour .... \unsure{L|E}orsque A est un anneau
local, on dira simplement par abus de langage que A est ``élémentaire de
multiplic\uline{i}té n'' si l'aug\var{\xout{m}}{}{typewriter strike-out}{ega25_25_161-v0014}mentation A $\hand{\to}$ A/$\underline{r}$(A) est élémentaire de
multiplicité \unsure{n}; noter que cela implique $n \hand{\geqslant} \unsure{2}$ , et si A est noethérien,
A est nécessairement non régulier. Si X est un préschéma, n un ent\unsure{i|e}\unsure{e|r}r

\mspage{\hand{161}}
